% Options for packages loaded elsewhere
\PassOptionsToPackage{unicode}{hyperref}
\PassOptionsToPackage{hyphens}{url}
\PassOptionsToPackage{dvipsnames,svgnames,x11names}{xcolor}
\documentclass[
  english,
  10pt,
  a4paper,
]{article}
\usepackage{xcolor}
\usepackage[top=22mm,bottom=20mm,left=18mm,right=18mm]{geometry}
\usepackage{amsmath,amssymb}
\setcounter{secnumdepth}{-\maxdimen} % remove section numbering
\usepackage{iftex}
\ifPDFTeX
  \usepackage[T1]{fontenc}
  \usepackage[utf8]{inputenc}
  \usepackage{textcomp} % provide euro and other symbols
\else % if luatex or xetex
  \usepackage{unicode-math} % this also loads fontspec
  \defaultfontfeatures{Scale=MatchLowercase}
  \defaultfontfeatures[\rmfamily]{Ligatures=TeX,Scale=1}
\fi
\usepackage{lmodern}
\ifPDFTeX\else
  % xetex/luatex font selection
\fi
% Use upquote if available, for straight quotes in verbatim environments
\IfFileExists{upquote.sty}{\usepackage{upquote}}{}
\IfFileExists{microtype.sty}{% use microtype if available
  \usepackage[]{microtype}
  \UseMicrotypeSet[protrusion]{basicmath} % disable protrusion for tt fonts
}{}
\makeatletter
\@ifundefined{KOMAClassName}{% if non-KOMA class
  \IfFileExists{parskip.sty}{%
    \usepackage{parskip}
  }{% else
    \setlength{\parindent}{0pt}
    \setlength{\parskip}{6pt plus 2pt minus 1pt}}
}{% if KOMA class
  \KOMAoptions{parskip=half}}
\makeatother
\usepackage{longtable,booktabs,array}
\usepackage{caption}
\captionsetup[table]{skip=6pt}
\newcounter{none} % for unnumbered tables
\usepackage{calc} % for calculating minipage widths
% Correct order of tables after \paragraph or \subparagraph
\usepackage{etoolbox}
\makeatletter
\patchcmd\longtable{\par}{\if@noskipsec\mbox{}\fi\par}{}{}
\makeatother
\ifLuaTeX
\usepackage[bidi=basic,shorthands=off]{babel}
\else
\usepackage[bidi=default,shorthands=off]{babel}
\fi
\ifLuaTeX
  \usepackage{selnolig} % disable illegal ligatures
\fi
\setlength{\emergencystretch}{3em} % prevent overfull lines
\providecommand{\tightlist}{%
  \setlength{\itemsep}{0pt}\setlength{\parskip}{0pt}}
% ---- Çılgın Kurye document theme: same look as the CSE 396 project proposal ----
% Compile with XeLaTeX (also works on Overleaf). Fonts: TeX Gyre Termes (Times), DejaVu Sans Mono.
\usepackage{fontspec}
\setmainfont{texgyretermes}[Extension=.otf, UprightFont=*-regular, BoldFont=*-bold, ItalicFont=*-italic, BoldItalicFont=*-bolditalic]
\setsansfont{texgyreheros}[Extension=.otf, UprightFont=*-regular, BoldFont=*-bold, ItalicFont=*-italic, BoldItalicFont=*-bolditalic]
\setmonofont{DejaVuSansMono}[Extension=.ttf, UprightFont=*, BoldFont=*-Bold, ItalicFont=*-Oblique, BoldItalicFont=*-BoldOblique, Scale=0.80]
\setmathfont{texgyretermes-math.otf}

\usepackage{xcolor}
\definecolor{ckblue}{HTML}{1F3A5F}
\definecolor{cklight}{HTML}{E8EEF5}
\definecolor{ckgrey}{HTML}{F2F2F2}
\definecolor{ckline}{HTML}{555555}
\definecolor{shadecolor}{HTML}{F4F5F7}

% Headings as in the proposal: serif, bold; subsections bold italic. Numbers are part of the heading text.
\setcounter{secnumdepth}{0}
\usepackage{titlesec}
\titleformat{\section}{\rmfamily\Large\bfseries}{}{0pt}{}
\titleformat{\subsection}{\rmfamily\large\bfseries\itshape}{}{0pt}{}
\titleformat{\subsubsection}{\rmfamily\normalsize\bfseries}{}{0pt}{}
\titlespacing*{\section}{0pt}{12pt}{4pt}
\titlespacing*{\subsection}{0pt}{8pt}{3pt}

% Header and footer (the first page has only the footer, like the proposal)
\usepackage{fancyhdr,lastpage}
\pagestyle{fancy}
\fancyhf{}
\fancyhead[L]{\footnotesize\itshape CSE 396 · Group 8 · Çılgın Kurye}
\fancyhead[R]{\footnotesize\itshape\ckmodule}
\fancyfoot[C]{\footnotesize Page \thepage\ of \pageref*{LastPage}}
\renewcommand{\headrulewidth}{0.4pt}
\setlength{\headheight}{14pt}
\fancypagestyle{ckfirst}{\fancyhf{}\renewcommand{\headrulewidth}{0pt}\fancyfoot[C]{\footnotesize Page \thepage\ of \pageref*{LastPage}}}

% Inline code: breaks only after . : ( , / _ and spaces, never inside a name
\usepackage[obeyspaces,spaces]{url}
\AtBeginDocument{%
  \renewcommand{\UrlBreaks}{\do\.\do\:\do\(\do\,\do\/\do\_\do\|}%
  \renewcommand{\UrlBigBreaks}{}%
  \renewcommand{\UrlNoBreaks}{}%
  \urlstyle{tt}}

% Tables: boxed grid as in the proposal, smaller text
\usepackage{etoolbox}
\AtBeginEnvironment{longtable}{\footnotesize\renewcommand{\arraystretch}{1.2}}
\setlength{\LTpre}{6pt}
\setlength{\LTpost}{8pt}
\setlength{\arrayrulewidth}{0.4pt}

% Code blocks: long lines wrap; diagrams stay in one piece
\usepackage{framed}
\usepackage{fvextra}
\RecustomVerbatimEnvironment{verbatim}{Verbatim}{fontsize=\footnotesize, breaklines}
\makeatletter
\@ifundefined{Shaded}{}{\renewenvironment{Shaded}{\begin{snugshade}}{\end{snugshade}}}
\@ifundefined{Highlighting}{}{\RecustomVerbatimEnvironment{Highlighting}{Verbatim}{commandchars=\\\{\}, fontsize=\footnotesize, breaklines, breakanywhere}}
\makeatother

% Figures drawn with TikZ
\usepackage{tikz}
\usetikzlibrary{arrows.meta,positioning,fit,calc,backgrounds,shapes.geometric}
\tikzset{
  ck/.style={font=\footnotesize, >={Stealth[length=5pt]}, line width=0.6pt},
  box/.style={draw=ckline, fill=white, rounded corners=2pt, align=center, inner sep=4pt, minimum height=7mm},
  mod/.style={box, fill=ckgrey},
  this/.style={box, fill=cklight, draw=ckblue, line width=0.9pt},
  group/.style={draw=ckline, dashed, rounded corners=3pt, inner sep=6pt},
  lbl/.style={font=\scriptsize, align=center, inner sep=1.5pt},
  msg/.style={->, draw=ckline},
}
\usepackage{caption}
\captionsetup[figure]{font=small, labelfont=bf, skip=4pt}
\newcommand{\ckfigure}[2]{\begin{figure}[htbp]\centering #1\caption{#2}\end{figure}}

\usepackage{enumitem}
\setlist{itemsep=1pt, topsep=3pt}
\setlength{\parskip}{4pt}
\setlength{\parindent}{0pt}
\raggedbottom
% ---- end of theme ----
\newcommand{\ckmodule}{M3 Physics and Collision}
\usepackage{bookmark}
\IfFileExists{xurl.sty}{\usepackage{xurl}}{} % add URL line breaks if available
\urlstyle{same}
% fallback for those not using the hyperref driver hyperxmp:
\makeatletter
\@ifundefined{xmpquote}{\newcommand{\xmpquote}[1]{#1}}{}
\makeatother
\hypersetup{
  pdftitle={M3 Physics and Collision},
  pdflang={en},
  colorlinks=true,
  linkcolor={black},
  filecolor={Maroon},
  citecolor={Blue},
  urlcolor={black},
  pdfcreator={LaTeX}}

\title{M3 Physics and Collision --- Module Documentation}
\author{}
\date{}


% ---- colours and drawings ----
% ---- colours ----
\definecolor{srv}{HTML}{2E6FB7}\colorlet{srvL}{srv!15}
\definecolor{cli}{HTML}{2E9E5B}\colorlet{cliL}{cli!15}
\definecolor{adm}{HTML}{E08A1E}\colorlet{admL}{adm!18}
\definecolor{brd}{HTML}{7B4FB3}\colorlet{brdL}{brd!15}
\definecolor{tcpc}{HTML}{1F5FA8}
\definecolor{udpc}{HTML}{D9531E}
\definecolor{obs}{HTML}{D7263D}
\definecolor{tA}{HTML}{00A6D6}
\definecolor{tB}{HTML}{3BB273}
\definecolor{tC}{HTML}{F2C14E}
\definecolor{tD}{HTML}{9B5DE5}

% ---- styles ----
\tikzset{
  nb/.style={rounded corners=3pt, line width=0.9pt, align=center, inner sep=4pt, minimum height=9mm, font=\footnotesize},
  sv/.style={nb, draw=srv, fill=srvL},
  cl/.style={nb, draw=cli, fill=cliL},
  ad/.style={nb, draw=adm, fill=admL},
  bd/.style={nb, draw=brd, fill=brdL},
  hwb/.style={nb, draw=black!50, fill=black!5},
  tcpa/.style={draw=tcpc, line width=1.5pt},
  udpa/.style={draw=udpc, line width=1.5pt, dash pattern=on 5pt off 2.5pt},
  inna/.style={draw=black!55, line width=1pt},
  al/.style={font=\scriptsize, fill=white, inner sep=1.2pt, rounded corners=1pt},
  grp/.style={draw=black!35, dashed, rounded corners=5pt, inner sep=7pt},
  glbl/.style={font=\scriptsize\bfseries, text=black!65},
  hl/.style={line width=2.6pt},
  st/.style={nb, rounded corners=8pt, minimum width=24mm},
  co/.style={draw=black!60, line width=0.5pt},
  cot/.style={font=\scriptsize, align=left, text=black!80},
}

% ---- system map; argument = module to highlight (M1 ... M8, or none) ----
\newcommand{\systemmap}[1]{%
\tikzset{M1/.style={},M2/.style={},M3/.style={},M4/.style={},M5/.style={},M6/.style={},M7/.style={},M8/.style={},none/.style={}}%
\tikzset{#1/.style={hl}}%
\begin{tikzpicture}[ck]
\node[cl, M5, minimum width=27mm] (m5) at (0,3) {\textbf{M5}\\Player phones};
\node[ad, M6, minimum width=27mm] (m6) at (0,0.4) {\textbf{M6}\\Admin panel};
\node[sv, M2, minimum width=27mm, minimum height=38mm] (m2) at (4.6,1.7) {\textbf{M2}\\Network\\and sync};
\node[sv, M1, minimum width=27mm, minimum height=38mm] (m1) at (8.4,1.7) {\textbf{M1}\\Game server\\and rules};
\node[sv, M3, minimum width=25mm] (m3) at (12,2.8) {\textbf{M3}\\Physics};
\node[sv, M4, minimum width=25mm] (m4) at (12,0.6) {\textbf{M4}\\Traffic and AI};
\node[grp, fit=(m2)(m1)(m3)(m4)] (pc) {};
\node[glbl, anchor=south west] at (pc.north west) {MAIN COMPUTER};
\node[bd, M7, minimum width=27mm] (m7) at (4.6,-2.7) {\textbf{M7}\\ESP32 software};
\node[bd, M8, minimum width=27mm] (m8) at (8.4,-2.7) {\textbf{M8}\\LED engine};
\node[hwb, minimum width=25mm] (led) at (12,-2.7) {344 LEDs};
\node[hwb, minimum width=27mm] (sm) at (4.6,-4.4) {Sensor and\\4 motors};
\node[grp, fit=(m7)(m8)(led)(sm)] (bb) {};
\node[glbl, anchor=north east] at (bb.south east) {LED CITY BOARD};
\draw[tcpa, <->] (m5.east) -- node[al]{WebSocket} (m5.east -| m2.west);
\draw[tcpa, <->] (m6.east) -- node[al]{WebSocket} (m6.east -| m2.west);
\draw[tcpa, <->] ([xshift=-6mm]m7.north) -- node[al]{TCP} ([xshift=-6mm]m2.south);
\draw[udpa, <->] ([xshift=6mm]m7.north) -- node[al]{UDP} ([xshift=6mm]m2.south);
\draw[inna, <->] (m2) -- (m1);
\draw[inna, <->] (m1.east |- m3) -- (m3);
\draw[inna, <->] (m1.east |- m4) -- (m4);
\draw[inna, ->] (m7) -- node[al]{state} (m8);
\draw[inna, ->] (m8) -- (led);
\draw[inna, <->] (m7) -- (sm);
\draw[tcpa] (-1.3,-6.0) -- ++(0.9,0) node[right, font=\scriptsize]{TCP: every message arrives, in order};
\draw[udpa] (4.0,-6.0) -- ++(0.9,0) node[right, font=\scriptsize]{UDP: fast, the newest message wins};
\draw[inna] (9.3,-6.0) -- ++(0.9,0) node[right, font=\scriptsize]{inside one program or chip};
\end{tikzpicture}}

% ---- LED board seen from above (6 x 6 roads, 1 unit = 10 cm) ----
\newcommand{\boardbase}{%
\fill[black!88, rounded corners=3pt] (-0.5,-0.5) rectangle (5.5,5.5);
\foreach \i in {0,...,4}{\foreach \j in {0,...,4}{%
  \fill[black!60] (\i+0.13,\j+0.13) rectangle (\i+0.87,\j+0.87);
  \fill[black!22] (\i+0.24,\j+0.24) rectangle (\i+0.52,\j+0.52);
  \fill[black!30] (\i+0.56,\j+0.30) rectangle (\i+0.78,\j+0.62);
  \fill[black!38] (\i+0.28,\j+0.58) rectangle (\i+0.66,\j+0.78);}}
\foreach \i in {0,...,5}{%
  \draw[black, line width=4.5pt] (\i,-0.05) -- (\i,5.05);
  \draw[black, line width=4.5pt] (-0.05,\i) -- (5.05,\i);}
\foreach \i in {0,...,5}{\foreach \k in {0,...,30}{%
  \fill[black!55] (\i,\k/6) circle (0.025);
  \fill[black!55] (\k/6,\i) circle (0.025);}}
\foreach \x/\y/\c in {0/4/tA, 4/4/tB, 0/0/tC, 4/0/tD}{%
  \fill[\c!30!black!60] (\x+0.13,\y+0.13) rectangle (\x+0.87,\y+0.87);
  \fill[\c] (\x+0.3,\y+0.3) rectangle (\x+0.7,\y+0.7);
  \fill[white] (\x+0.4,\y+0.5) circle (0.05); \fill[white] (\x+0.6,\y+0.5) circle (0.05);}}
\newcommand{\nmark}[3]{\node[circle, fill=yellow!85!orange, draw=black, line width=0.4pt, inner sep=0.6pt, minimum size=3.6mm, font=\tiny\bfseries] at (#1,#2) {#3};}
\newcommand{\bike}[4]{%
\fill[#3!35] ($(#1,#2)+(#4:0.40)$) circle (0.05);
\fill[#3!65] ($(#1,#2)+(#4:0.21)$) circle (0.075);
\fill[#3, draw=white, line width=0.4pt] (#1,#2) circle (0.12);}
\newcommand{\dpoint}[3]{\draw[#3!50, line width=0.8pt] (#1,#2) circle (0.26); \draw[#3, line width=1.4pt] (#1,#2) circle (0.16); \fill[#3] (#1,#2) circle (0.07);}
\newcommand{\heavy}[2]{\draw[obs!60!black, line width=4pt] (#1) -- (#2); \draw[obs, line width=4pt, dash pattern=on 3pt off 5pt] (#1) -- (#2);}
\newcommand{\roadwork}[2]{\draw[obs, line width=4pt] (#1) -- (#2);}
\newcommand{\accident}[2]{\draw[obs, line width=4pt] (#1) -- (#2); \draw[white, line width=4pt, dash pattern=on 2.5pt off 3.5pt] (#1) -- (#2);}
\newcommand{\puddle}[2]{\draw[cyan!35, line width=4pt] (#1) -- (#2); \draw[white, line width=1.6pt, dash pattern=on 1pt off 2.5pt] (#1) -- (#2);}
\newcommand{\findme}[2]{\fill[white, opacity=0.25] (#1,#2) circle (0.36); \fill[white, opacity=0.5] (#1,#2) circle (0.24); \fill[white] (#1,#2) circle (0.13);}

% ---- phone game screen (7 x 3.4 units) ----
\newcommand{\phonemock}{%
\fill[black!85, rounded corners=9pt] (-0.3,-0.3) rectangle (7.3,3.7);
\fill[cyan!20] (0,1.7) rectangle (7,3.4);
\fill[cli!14] (0,0) rectangle (7,1.7);
\filldraw[fill=srv!30, draw=srv!70!black] (0.2,1.5) rectangle (1.0,3.0);
\filldraw[fill=adm!35, draw=adm!70!black] (1.1,1.6) rectangle (1.8,2.4);
\filldraw[fill=brd!30, draw=brd!70!black] (1.9,1.65) -- (2.6,1.65) -- (2.6,2.2) -- (2.25,2.5) -- (1.9,2.2) -- cycle;
\filldraw[fill=adm!35, draw=adm!70!black] (4.5,1.65) rectangle (5.1,2.3);
\filldraw[fill=srv!30, draw=srv!70!black] (5.2,1.55) rectangle (6.0,3.1);
\filldraw[fill=cli!35, draw=cli!70!black] (6.1,1.6) rectangle (6.85,2.6);
\fill[brown!60!black] (4.17,1.6) rectangle (4.23,1.8); \fill[cli!70!black] (4.2,1.95) circle (0.18);
\fill[black!55] (1.8,0) -- (5.2,0) -- (3.75,1.7) -- (3.25,1.7) -- cycle;
\draw[white, line width=0.8pt, dash pattern=on 4pt off 4pt] (3.5,0) -- (3.5,1.7);
\filldraw[fill=black!15, draw=black!70] (3.62,1.12) rectangle (3.88,1.32);
\fill[obs] (3.05,1.22) -- (3.15,1.48) -- (3.25,1.22) -- cycle;
\fill[black!85] (3.43,0.08) rectangle (3.57,0.3);
\fill[tA] (3.3,0.28) rectangle (3.7,0.6);
\fill[brown!65] (3.34,0.6) rectangle (3.66,0.82);
\fill[tA!60!black] (3.5,0.95) circle (0.13);
\node[fill=black!70, text=white, rounded corners=2pt, font=\scriptsize\bfseries, inner sep=2pt] at (3.5,3.15) {2:41};
\foreach \c/\s/\x in {tA/120/0.35,tB/95/0.95,tC/140/1.55,tD/80/2.15}
  \node[fill=\c, text=white, font=\tiny\bfseries, rounded corners=1pt, inner sep=1.5pt, minimum width=5mm] at (\x,3.15) {\s};
\fill[white, opacity=0.9] (6.45,2.85) circle (0.38); \draw[black!60] (6.45,2.85) circle (0.38);
\draw[-{Stealth[length=5pt]}, line width=1.8pt, obs] (6.45,2.85) ++(230:0.2) -- ++(50:0.5);
\fill[white, opacity=0.85] (0.55,0.5) circle (0.32); \fill[black!70] (0.43,0.5) -- (0.65,0.63) -- (0.65,0.37) -- cycle;
\fill[white, opacity=0.85] (1.4,0.5) circle (0.32); \fill[black!70] (1.52,0.5) -- (1.3,0.63) -- (1.3,0.37) -- cycle;
\node[fill=obs!80, text=white, rounded corners=3pt, font=\tiny\bfseries, minimum width=8mm, minimum height=6mm] at (5.6,0.5) {Brake};
\node[fill=cli, text=white, rounded corners=3pt, font=\tiny\bfseries, minimum width=8mm, minimum height=8mm] at (6.5,0.6) {Gas};
\node[fill=tC, text=black, rounded corners=3pt, font=\tiny\bfseries, minimum width=9mm, minimum height=5mm] at (6.15,1.45) {Deliver};
\node[fill=white, text=black, rounded corners=3pt, font=\tiny\bfseries, minimum width=9mm, minimum height=5mm] at (0.9,1.45) {Find me};}

% ---- side cut of the board (14 units wide) ----
\newcommand{\boardside}{%
\fill[brown!20] (-0.6,-0.35) rectangle (14.6,0);
\foreach \x/\h/\c in {1.0/0.7/black!18, 2.2/1.1/black!22, 4.6/1.9/black!25, 5.8/2.6/black!28, 7.4/2.2/black!25, 8.6/1.6/black!22, 10.9/1.0/black!18}
  \filldraw[fill=\c, draw=black!45] (\x,1.5) rectangle (\x+0.9,1.5+\h);
\filldraw[fill=tB!35, draw=tB!70!black] (12.4,1.5) rectangle (13.4,2.3);
\fill[black!70, rounded corners=1pt] (0.5,0) rectangle (1.1,0.9);
\fill[black!70, rounded corners=1pt] (12.9,0) rectangle (13.5,0.9);
\filldraw[fill=brown!45, draw=brown!70!black] (0,0.9) rectangle (14,1.3);
\fill[black!80] (0,1.3) rectangle (14,1.38);
\foreach \x in {0.3,0.9,...,13.8} \fill[black!35] (\x,1.42) circle (0.06);
\foreach \x/\c in {3.3/tA, 3.9/tA, 6.9/obs, 7.5/obs, 10.5/tC} \fill[\c] (\x,1.42) circle (0.07);
\fill[white, opacity=0.55] (0,1.38) rectangle (14,1.52);
\filldraw[fill=brd!40, draw=brd] (2.3,0.35) rectangle (4.3,0.85);
\node[font=\tiny] at (3.3,0.6) {ESP32};
\filldraw[fill=brd, draw=brd!60!black] (6.75,0.62) rectangle (7.25,0.88);
\filldraw[fill=adm!60, draw=adm!70!black] (1.8,0.6) circle (0.24);
\filldraw[fill=adm!60, draw=adm!70!black] (12.2,0.6) circle (0.24);
\filldraw[fill=black!15, draw=black!50] (9.3,0.35) rectangle (10.9,0.85);
\node[font=\tiny] at (10.1,0.6) {Power};}

\begin{document}

\thispagestyle{ckfirst}
\begin{center}
{\bfseries GEBZE TECHNICAL UNIVERSITY\par}
{Computer Engineering Department\par}
{CSE 396 Module Documentation\par}
\vspace{8pt}
{\LARGE\bfseries Çılgın Kurye\par}
\vspace{3pt}
{\large\itshape Module M3: Physics and Collision\par}
\vspace{6pt}
\textbf{Instructor:} Erkan Zergeroğlu\hspace{1.4em}\textbf{Company:} Board\&Byte\hspace{1.4em}\textbf{Group:} Group-8\hspace{1.4em}\textbf{Date:} 6 November 2026
\end{center}
\begin{center}\footnotesize\renewcommand{\arraystretch}{1.15}
\begin{tabular}{|p{0.27\linewidth}|p{0.15\linewidth}|p{0.27\linewidth}|p{0.15\linewidth}|}\hline
\textbf{Name} & \textbf{Student No.} & \textbf{Name} & \textbf{Student No.} \\ \hline
Abdullatif Yılmaz & 230104004031 & Hasan Alperen Elmas & 1801042604 \\ \hline
Batuhan Yanar & 220104004059 & Selim Barut & 220104004046 \\ \hline
\end{tabular}
\end{center}
\vspace{2pt}



% ---- table and list helpers ----
\newcolumntype{L}[1]{>{\raggedright\arraybackslash}p{\dimexpr#1\linewidth-2\tabcolsep-\arrayrulewidth\relax}}
\newcolumntype{C}[1]{>{\centering\arraybackslash}p{\dimexpr#1\linewidth-2\tabcolsep-\arrayrulewidth\relax}}
\newcommand{\ckdone}{\item[$\boxtimes$]}
\newcommand{\ckopen}{\item[$\square$]}
\newcommand{\mdot}{$\bullet$}
\newcounter{msec}
\newcounter{chap}
% ---- end of helpers ----
\newcommand{\msection}[1]{\stepcounter{msec}\section{\arabic{msec}. #1}}

\msection{Purpose}

M3 moves the motorcycles. It turns throttle, brake and steering into speed, direction and position.
It finds collisions, decides when a motorcycle falls, and adds the effect of obstacles and
earthquakes to the motion. The same code runs in two places: on the server, where its result is
the real one, and on the phone, where it guesses the player's own motorcycle before the server
answers.

\begin{longtable}[]{|L{0.18}|L{0.82}|}
\hline Runs on & Main computer inside the game server; phone browser for prediction only \\ \hline
Technology & TypeScript library; our own simple 2D arcade physics, no ready physics engine; circle and box collisions on a uniform grid; seeded random numbers \\ \hline
Requirements & R3 R4 R6 R7 R11 \\ \hline
\end{longtable}

\begin{figure}[htbp]\centering
\systemmap{M3}
\caption{Where M3 sits in the system. M1 calls M3 in every tick; M5 uses a copy of M3 on the phone.}
\end{figure}

\msection{Requirements covered}

\begin{longtable}[]{|L{0.07}|L{0.25}|L{0.68}|}
\hline \textbf{ID} & \textbf{Requirement} & \textbf{What M3 does for it} \\ \hline \endhead
R3 & Driving and controls & Throttle, brake, left and right; no reverse; U-turn by steering in about a 3 m circle at low speed \\ \hline
R4 & Collisions and falls & Collisions with buildings, barriers, NPC cars and motorcycles; a fast hit makes the motorcycle fall \\ \hline
R6 & Obstacles & Road work and accidents are solid; an accident has a slow zone; a puddle makes the motorcycle slide \\ \hline
R7 & Earthquake & Side push and less grip inside the earthquake radius that M1 sends \\ \hline
R11 & Single authority & The code gives the same result for the same input; the server result always wins; the phone only predicts \\ \hline
\end{longtable}

\msection{Inputs and outputs}

\begin{longtable}[]{|L{0.08}|L{0.20}|L{0.18}|L{0.54}|}
\hline \textbf{Dir.} & \textbf{Data} & \textbf{From / to} & \textbf{Content} \\ \hline \endhead
In & City plan & Shared file & Roads, lanes, blocks and buildings as fixed colliders \\ \hline
In & Controls & M1 & Every tick: throttle, brake, steer of humans and bots; a speed factor for each player \\ \hline
In & Obstacles, NPC cars & M1, cars from M4 & Obstacle type, road, lane, position; boxes of the NPC cars \\ \hline
In & Earthquake & M1 & Centre, strength, radius; the players M1 decided to drop \\ \hline
In & Parameters & M6, through M1 & Physics group: top speed, grip, fall speeds, puddle and earthquake effects \\ \hline
Out & Poses & M1 & Position, direction, speed, lean, sliding, fallen \\ \hline
Out & Events & M1 and M4 & Collisions (with what, how fast) and falls \\ \hline
Out & Predicted pose & M5 & The player's own motorcycle before the server answers \\ \hline
\end{longtable}

\msection{Main components}

\begin{longtable}[]{|L{0.20}|L{0.52}|L{0.28}|}
\hline \textbf{Component} & \textbf{What it does} & \textbf{Related modules} \\ \hline \endhead
Motorcycle model & Speed, steering, turning, grip and sliding & --- \\ \hline
Colliders & Buildings and the outer wall from the city plan; obstacles as boxes, slow zones and slippery areas & M1 \\ \hline
Collision solver and fall judge & Solves contacts in a fixed order and decides falls from the hit speed & M1, M4 \\ \hline
Earthquake effect & Push and less grip inside the radius & M1 \\ \hline
Predictor & Prediction and correction on the phone & M5 \\ \hline
\end{longtable}

\msection{How it works}

\textbf{One physics step.} The server runs 30 ticks per second, and M3 runs 2 sub-steps in each
tick, so one sub-step is about 16 ms. Small steps keep fast motorcycles from passing through thin
walls. Each sub-step does five things in a fixed order, as the figure shows. A fixed order and
seeded random numbers mean that the same inputs always give the same result on the server and on
the phone.

\begin{figure}[htbp]\centering
\begin{tikzpicture}[ck, node distance=4mm]
\node[cl, minimum width=19mm] (a) {Speed and\\steering};
\node[sv, minimum width=19mm, right=of a, fill=tC!25, draw=tC!70!black] (b) {Grip and\\sliding};
\node[sv, minimum width=19mm, right=of b] (c) {Walls and\\obstacles};
\node[sv, minimum width=19mm, right=of c, fill=tD!18, draw=tD!80!black] (d) {Cars and\\motorcycles};
\node[nb, draw=obs, fill=obs!10, minimum width=19mm, right=of d] (e) {Fall\\check};
\foreach \u/\v in {a/b,b/c,c/d,d/e} \draw[inna,->] (\u) -- (\v);
\draw[inna,->, srv] (e.south) -- ++(0,-0.55) -| node[al, pos=0.25]{2 sub-steps in each tick} (a.south);
\end{tikzpicture}
\caption{The five parts of one physics sub-step.}
\end{figure}

\textbf{Driving.} The throttle adds speed up to the top speed, and the brake and air drag take speed
away. There is no reverse. The steering changes the direction. At low speed the motorcycle can turn
in a circle of about 3 metres, so a U-turn is easy. At high speed the turn is wider. The model also
gives a lean angle, which M5 uses to tilt the 3D motorcycle in a turn.

\textbf{Grip and sliding.} In a turn the motorcycle needs a side force. M3 compares the size of this
force with the grip limit. If the force is below the limit, the motorcycle follows the turn. If it
is above the limit, the motorcycle slides out of the turn. We compare only the size, so a left turn
and a right turn act the same. On a water puddle the grip limit is very low, so the motorcycle
slides but does not slow down. Inside an earthquake radius the grip is lower too, and there is a
small push to the side.

\begin{figure}[htbp]\centering
\begin{tikzpicture}[ck]
\begin{scope}
\fill[black!65] (-0.5,-0.5) rectangle (4.5,3.2);
\draw[white, dash pattern=on 5pt off 5pt] (2,-0.5) -- (2,3.2);
\draw[cli, line width=1.6pt, ->] (1,-0.2) .. controls (1,1.5) and (1.6,2.4) .. (3.6,2.6);
\fill[tA] (1,0.6) circle (0.15);
\node[font=\scriptsize, text=white, align=center] at (3.3,0.6) {grip is\\enough};
\node[font=\scriptsize\bfseries] at (2,-0.85) {Normal turn};
\end{scope}
\begin{scope}[shift={(5.6,0)}]
\fill[black!65] (-0.5,-0.5) rectangle (4.5,3.2);
\draw[white, dash pattern=on 5pt off 5pt] (2,-0.5) -- (2,3.2);
\draw[cli!60, line width=1.2pt, dashed] (1,-0.2) .. controls (1,1.5) and (1.6,2.4) .. (3.6,2.6);
\draw[obs, line width=1.6pt, ->] (1,-0.2) .. controls (1,1.4) and (1.1,2.2) .. (1.2,3.1);
\fill[tA] (1,0.6) circle (0.15);
\node[font=\scriptsize, text=white, align=center] at (3.3,0.6) {too fast,\\slides out};
\node[font=\scriptsize\bfseries] at (2,-0.85) {Too fast in a turn};
\end{scope}
\begin{scope}[shift={(11.2,0)}]
\fill[black!65] (-0.5,-0.5) rectangle (4.5,3.2);
\draw[white, dash pattern=on 5pt off 5pt] (2,-0.5) -- (2,3.2);
\fill[cyan!35] (1.0,1.2) ellipse (0.9 and 0.5);
\draw[white, line width=0.8pt, dash pattern=on 1pt off 2pt] (1.0,1.2) ellipse (0.6 and 0.3);
\draw[obs, line width=1.6pt, ->] (1,-0.2) -- (1,0.9) .. controls (1.1,1.5) and (1.6,1.8) .. (2.6,2.9);
\fill[tA] (1,0.4) circle (0.15);
\node[font=\scriptsize, text=white, align=center] at (3.3,0.6) {puddle,\\low grip};
\node[font=\scriptsize\bfseries] at (2,-0.85) {On a puddle};
\end{scope}
\end{tikzpicture}
\caption{Three cases of grip, seen from above. Green is the wanted path; red is the real path.}
\end{figure}

\textbf{Collisions and falls.} Buildings, the outer wall and the road-work barriers come from the city
plan and the obstacle list as boxes. NPC cars are boxes, and motorcycles are circles. To find
collisions fast, M3 puts every object into a uniform grid, so it checks only objects in the same or
next cell. A contact pushes the objects apart. If the hit speed is above the fall speed that the
admin set, the motorcycle falls. When two motorcycles hit each other and their relative speed is
above the limit, both of them fall. A fallen motorcycle is not solid, so it does not block the road
for others. M1 gives the penalty time and later asks M3 to stand the motorcycle up again. An
accident is solid in one lane and has a slow zone around it.

\textbf{Prediction on the phone.} On the phone, the player must see the own motorcycle move at once,
not 100 ms later. So M5 runs a copy of M3 for this one motorcycle. Each control gets a number and is
kept in a list. When a server state arrives, the phone takes the real position from the server,
throws away the controls that the server has already used, and runs the remaining controls again.
If the new result is close to what the player sees, the difference is blended away in a short
time. If it is far, for example after a crash, the motorcycle jumps to the right place.

\begin{figure}[htbp]\centering
\begin{tikzpicture}[ck]
\draw[black!40, ->] (0,0) -- (13,0) node[right, font=\scriptsize]{time};
\foreach \x/\l in {1/1,3/2,5/3,7/4,9/5} {\fill[cli] (\x,0.9) circle (0.09); \node[font=\tiny] at (\x,1.2) {control \l};}
\node[font=\scriptsize\bfseries, text=cli, anchor=east] at (-0.1,0.9) {Phone};
\node[font=\scriptsize\bfseries, text=srv, anchor=east] at (-0.1,-0.9) {Server};
\draw[tcpa,->] (1,0.8) -- (4.0,-0.8);
\draw[tcpa,->] (3,0.8) -- (6.0,-0.8);
\node[sv, minimum height=6mm, font=\tiny] (s) at (6.6,-0.9) {state with control 2};
\draw[tcpa,->] (6.6,-0.6) -- (9.6,0.75);
\node[cl, minimum height=6mm, font=\tiny, align=center] at (11.6,0.9) {take server position,\\run controls 3 to 5 again};
\fill[cli!25, rounded corners=2pt] (4.6,0.55) rectangle (9.4,0.65);
\node[font=\tiny, text=cli!70!black] at (7,0.4) {controls 3 to 5 are not yet known by the server};
\end{tikzpicture}
\caption{Prediction and correction for the player's own motorcycle.}
\end{figure}

\msection{Testing and demonstration}

M3 needs no other module. A test makes a world from the city plan and gives scripted controls. We
test a straight run and braking, a U-turn, slides to the left and to the right, wall and car
crashes, the earthquake radius, the same result for the same inputs, and prediction with delay and
loss. In the demonstration one motorcycle drives in the city with the keyboard. The admin changes
grip and top speed, and we see the motorcycle change at once.

\end{document}
